% TimblServer 1.20 manual

\documentclass{report}

\usepackage{a4wide}
\usepackage{palatino}
\usepackage{url}

% allow long URLs to break across lines
\urldef{\servermanualurl}\url{https://github.com/LanguageMachines/timblserver/blob/master/docs/TimblServer_1.20_Manual.pdf}
\urldef{\timblmanualurl}\url{https://github.com/LanguageMachines/timbl/blob/master/docs/Timbl_7.0_Manual.pdf}

\author{Ko van der Sloot \\
	Centre for Language Studies \\
        Radboud University Nijmegen \\ \\
       URL: https://languagemachines.github.io/timbl\thanks{This document is available from
	\servermanualurl. All rights reserved
	Induction of Linguistic Knowledge, Tilburg University and
        CLiPS, University of Antwerp.}}

\title{{\huge TimblServer: Tilburg Memory-Based Learner Server} \\ \vspace*{0.5cm}
{\bf version 1.20} \\ \vspace*{0.5cm}{\huge Manual}}

%better paragraph indentation
\parindent 0pt
\parskip 9pt

\begin{document}

\pagenumbering{roman}

\maketitle

\tableofcontents

\chapter*{Preface}

This reference guide describes the server interface to TiMBL, the
Tilburg Memory-Based Learner software package. To be able to use
TimblServer, a working installation of TiMBL (version 7.0 or later)
should be present on your system. For information on the installation
and use of TiMBL, see the TiMBL Reference Guide \cite{Daelemans+26}
and the TiMBL webpage\footnote{\url{https://languagemachines.github.io/timbl}}.

TimblServer wraps one or more TiMBL classifiers (instance bases, or
henceforth {\em bases}) in a
network server, so that trained classifiers can be kept in memory and
be queried from multiple simultaneous client sessions, locally or over
the network. Since version 1.13 the server speaks three protocols: a
line-based TCP protocol, an HTTP protocol with XML output, and a JSON
protocol.

Although this software has been tested for a long time in our research
groups, it may still contain bugs and inconsistencies in some
places. We would appreciate it if you would file bug reports, ideas
about enhancements of the software and the manual, and any other
comments you might have, as issues at
\url{https://github.com/LanguageMachines/timblserver/issues}.

This reference guide is structured as follows. In
Chapter~\ref{license} you can find the terms of the license according
to which you are allowed to use TimblServer. The subsequent chapter
gives some instructions on how to install the TimblServer package on
your computer. Chapter~\ref{changes} lists the changes that have taken
place up to the current version. Next, Chapter~\ref{tutorial} offers a
quick-start tutorial for readers who want to get to work with
TimblServer right away. Chapter~\ref{reference} describes the command
line options and the configuration file, and
Chapter~\ref{serverformat} documents the three server protocols.
Chapter~\ref{client} describes the {\tt timblclient} test program.

\chapter{GNU General Public License}
\label{license}
\pagenumbering{arabic}

TimblServer is free software; you can redistribute it and/or modify it
under the terms of the GNU General Public License as published by the
Free Software Foundation; either version 3 of the License, or (at your
option) any later version.

TimblServer is distributed in the hope that it will be useful, but
WITHOUT ANY WARRANTY; without even the implied warranty of
MERCHANTABILITY or FITNESS FOR A PARTICULAR PURPOSE.  See the GNU
General Public License for more details.

You should have received a copy of the GNU General Public License
along with TimblServer.  If not, see
$<$http://www.gnu.org/licenses/$>$.

In publication of research that makes use of TimblServer, a citation
should be given of: {\em ``Walter Daelemans, Jakub Zavrel, Ko van der
  Sloot, and Antal van den Bosch (2026). TiMBL: Tilburg Memory Based
  Learner, version 7.0, Reference Guide. Available from} \\
{\tt https://github.com/LanguageMachines/timbl/blob/master/docs/}\\
{\tt Timbl\_7.0\_Manual.pdf}''.

\pagestyle{headings}

\chapter{Installation}

The TimblServer source code is hosted on GitHub:

{\tt https://github.com/LanguageMachines/timblserver}

Stable releases can be downloaded as tarballs from

{\tt https://github.com/LanguageMachines/timblserver/releases}

and the development version can be obtained with

{\tt > git clone https://github.com/LanguageMachines/timblserver}

The package contains the complete source code (C++) for the
TimblServer program, the license, and documentation. The installation
should be relatively straightforward on most UNIX systems.

Before installing TimblServer, please make sure the following
dependencies are installed on your system:

\begin{itemize}
\item {\tt ticcutils}\footnote{\url{https://github.com/LanguageMachines/ticcutils}}
\item {\tt timbl} (version 7.0 or later)\footnote{\url{https://github.com/LanguageMachines/timbl}}
\item {\tt libxml2}
\item {\tt pkg-config}
\end{itemize}

Building requires a C++17 compiler (e.g.\ GCC 7 or later, or a recent
Clang). If in doubt about the installed TiMBL version, check the
output of the command {\tt timbl -V}.

To compile and install from a cloned repository ({\tt >} is the
command line prompt):

{\tt > cd timblserver} \\
{\tt > bash bootstrap.sh} \\
{\tt > ./configure --prefix=<location\_to\_install>} \\
{\tt > make} \\
{\tt > make install}

When building from a release tarball, the {\tt bootstrap.sh} step can
be skipped. If you do not use the {\tt --prefix} option, TimblServer
will try to install itself in the directory {\tt /usr/local/}. If you
do not have {\tt root} access you can specify a different installation
location such as {\tt \$HOME/install}.

If the process was completed successfully, you should now have
executable files named {\tt timblserver} and {\tt timblclient} in the
directory {\tt <location\_to\_install>/bin}, and a library {\tt
  libtimblserver} in the directory {\tt <location\_to\_install>/lib}.

Alternatively, TimblServer is packaged by several Linux distributions,
and it is included in our LaMachine software
distribution\footnote{\url{https://proycon.github.io/LaMachine}}.

The address for problems with the installation, bug reports, comments
and questions is our issue tracker at
\url{https://github.com/LanguageMachines/timblserver/issues}.

\chapter{Changes}
\label{changes}

\section{From version 1.1 to 1.20}

The changes most relevant to users are summarised below; the {\tt
  NEWS} and {\tt ChangeLog} files in the source repository give the
complete per-release history.

\begin{itemize}
\item The executables were renamed from {\tt TimblServer} and {\tt
  TimblClient} to lowercase {\tt timblserver} and {\tt timblclient},
  and the library to {\tt libtimblserver} (version 1.2).
\item Next to the original line-based TCP protocol, an HTTP protocol
  was added, which answers {\tt GET} requests with XML; the {\tt
  protocol} attribute in the configuration file (marked ``reserved for
  future usage'' in the 1.1 manual) selects it. See
  Section~\ref{httpproto}.
\item JSON support was added in version 1.13 (2019). With {\tt
  protocol=json} the server communicates in newline-delimited JSON
  objects instead of the line-based TCP protocol. The JSON syntax
  might still change in future versions, so handle with care. See
  Section~\ref{jsonproto}.
\item The generic server infrastructure was moved to the {\tt
  ticcutils} package (versions 1.6--1.10), on which TimblServer now
  depends, and the source repository moved to GitHub (version 1.10).
\item The long options {\tt --version} and {\tt --help} are accepted
  alongside {\tt -V} and {\tt -h} (version 1.11).
\item The dependencies on Boost and libtar were removed; TimblServer
  now depends on {\tt ticcutils}, {\tt timbl}, {\tt libicu} and {\tt
  libxml2} (versions 1.15--1.17).
\item Building TimblServer now requires a C++17 compiler (version
  1.19).
\end{itemize}

\section{version 1.1}

This is the initial release of a separate TimblServer.  Compared with
the version included in Timbl (up to version 6.2) the following is
changed:

\begin{itemize}
\item Some bugs are fixed ({\tt -C <n>} didn't work as expected)
\item a {\tt --config} option is added to read a configuration file. This
  file can be used to specify a whole range of separated and unrelated
  TiMBL experiments all to run on the same TCP port.
\end{itemize}

\chapter{Quick-start Tutorial}
\label{tutorial}

The use of TimblServer subsumes a familiarity with TiMBL.  The
TimblServer command line interface is the same as that of {\tt timbl}
with a few additions, exemplified below.

If you want to run the examples and demos from this manual, you should
act as follows. We will use the same examples as TiMBL, so copy the
files {\tt dimin.train} and {\tt dimin.test} from the TiMBL
distribution\footnote{In the {\tt demos} directory of the TiMBL source
  tree.} to some working location. From the directory you copied the
files to, issue the following command:

{\tt > timblserver -f dimin.train -S 7000}

This will start a TimblServer on port 7000 (if that port is occupied
you might need to try another number).

The server can be used (from the same machine) over Telnet like this:

{\tt > telnet localhost 7000 }

(or, equivalently, {\tt nc localhost 7000}). The server responds:

{\tt Welcome to the Timbl server.}

Try to classify an instance:

{\tt classify =,=,=,=,=,=,=,=,+,p,e,=,? }

The server should respond with:

{\tt CATEGORY \{T\} }

Type {\tt exit} to close the session. The server keeps running, and
must be stopped by killing the {\tt timblserver} process.

\chapter{Software usage and options}
\label{reference}

\section{Command line options}
\label{commandline}

The user interacts with TimblServer through the use of command line
arguments.  When you have installed TimblServer successfully, and you
type {\tt timblserver} at the command line without any further
arguments, it will print an overview of the most basic command line
options:

{\footnotesize
\begin{verbatim}
TiMBL Server 1.20 (c) CLTS/ILK/CLIPS 1998 - 2026.
Tilburg Memory Based Learner
Centre for Language and Speech Technology, Radboud University
Induction of Linguistic Knowledge Research Group, Tilburg University
CLiPS Computational Linguistics Group, University of Antwerp
based on timbl 7.0

usage:  timblserver --config=config-file
        timblserver -f data-file {-S socket} {-C num}
or see: timblserver -h
        for more options
\end{verbatim}
}

\subsection{Server options}

\begin{description}

\item {\tt --config=<file>} or {\tt -c <file>} : Start a server with a
  range of TiMBL experiments as specified in the configfile.

This is the {\bf preferred} way to use TimblServer. See
Section~\ref{configfile} for details about the configfile.

\item {\tt --protocol=[tcp|http|json]} : select the protocol the
  server speaks (default {\tt tcp}). See Chapter~\ref{serverformat}.
  The protocol can also be set in the configfile.

\item {\tt --pidfile=<file>} : store the pid of the server in 'file'.
  A full path should be given, otherwise the file will end up relative
  to '/'.

\item {\tt --logfile=<file>} : log server messages to 'file'.  A full
  path should be given, otherwise the file will end up relative to
  '/'.

\item {\tt --daemonize=[yes|no] } : run the server as a daemon. The
  default is {\em yes}.

\item {\tt -S <portnumber>} : Starts a TiMBL server listening on the
  specified port number of the localhost.  This option is deprecated
  and only available to be backward compatible with Timbl 6.2 and
  earlier versions; the {\tt --config} variant is preferred.

\item {\tt -C <number>} : limit the number of parallel connections to
  {\tt <number>}. The default is 10. Deprecated, like {\tt -S}.

\end{description}

All other options are passed on to the embedded TiMBL experiment(s);
they are documented in the TiMBL Reference Guide \cite{Daelemans+26}
and in {\tt timbl}(1). The long options {\tt --version} and {\tt
  --help} are accepted as aliases for {\tt -V} and {\tt -h}.

\subsection{Configuration file}
\label{configfile}

The configfile contains lines of the general format {\tt
  attribute=value}. Valid attributes are: {\tt port} to set the TCP
port, {\tt maxconn} to set the number of parallel connections, {\tt
  protocol} to select the {\tt tcp}, {\tt http} or {\tt json} protocol
(default {\tt tcp}), and the equivalents of the command line options
above: {\tt pidfile}, {\tt logfile}, {\tt daemonize}. The attribute
{\tt debug} can be set to get (very) verbose logging.

Besides these attribute--value pairs, the configfile contains any
number of lines describing TiMBL experiments to be served. The lines
take the form of {\tt basename="timbl training options"}.  {\tt
  basename} is a name you can choose to identify the experiment. The
{\tt basename} is needed when you call the server to classify an
instance.

{\tt "timbl training options"} is a string of options which are used
to train Timbl. For example:

\begin{verbatim}
port=7000
maxconn=10
protocol=tcp
[[experiments]]
dimin0="-aIB1 +vdi+db+n -f dimin.train"
dimin1="-aIGTREE +vdi+db +D -f dimin.train"
dimin2="-aTRIBL +vdi+db+n -q3 -f dimin.train"
\end{verbatim}

If you start TimblServer with this configfile, it will start on port
7000, with at most 10 parallel sessions. The server will train 3
different Timbl instances, identified by the names {\tt dimin0}, {\tt
  dimin1} and {\tt dimin2}.

The optional section header {\tt [[experiments]]} separates the
experiment definitions from the global server settings. For backward
compatibility it may be omitted, in which case every
attribute--value pair that is not a recognised server setting is
interpreted as an experiment definition.

Instead of training from a data file with {\tt -f}, an experiment can
also load a previously stored instance base with {\tt -i}, which makes
the server start up much faster; optional weights ({\tt -w}),
probability arrays ({\tt -u}) and value-difference matrices ({\tt
  --matrixin}) files are read as well, exactly as in standalone {\tt
  timbl}.

\chapter{Server protocols}
\label{serverformat}

When the TimblServer is running, it is waiting for input on the
specified port number. When a client connects on this port number, the
server starts a separate thread, processing any given commands, such
as to classify a new example. Which protocol the server speaks on that
port is determined by the {\tt protocol} setting: {\tt tcp} (the
default, Section~\ref{tcpproto}), {\tt http}
(Section~\ref{httpproto}), or {\tt json} (Section~\ref{jsonproto}).
One server instance speaks one protocol.

\section{The TCP protocol}
\label{tcpproto}

The TCP protocol is a simple line-based protocol, suitable for use
with {\tt telnet} or {\tt nc}, and the protocol spoken by {\tt
  timblclient} (Chapter~\ref{client}). After accepting the connection,
the server first sends a welcome message to the client:

{\tt Welcome to the Timbl server.}

When the server hosts more than one base, this is followed by a line
listing the available bases:

{\tt available bases: dimin0 dimin1 dimin2 }

After this, the server waits for client-side requests.  The client can
now issue five types of commands: {\tt base}, {\tt classify}, {\tt
  set} (options), {\tt query} (status), and {\tt exit}. The type of
command is specified by the first string of the request line,
which can be abbreviated to any prefix of the command, up to one
letter (i.e.~b, c, s, q, e); commands are case-insensitive. The
command is followed by whitespace and the remainder of the command as
described below. Lines starting with {\tt \#} are skipped.

\begin{description}
\item {\tt base basename}\\
  tell the server to which base the next command(s) refer. The
  basename must be one of the basenames specified in the configfile.
  If no configfile was specified, i.e. TimblServer was
  started with the old-school {\tt -S} option, then the basename is
  ``default''. It is not necessary to specify this default name when
  testing (therefore preserving backward compatibility with older
  Timbl versions).

\item {\tt classify testcase}\\
  testcase is a pattern of features (must have the same number of
  features as the training set) followed by a category
  string. E.g.: {\tt small,long,1,??.}\\
  Depending on the current settings of the server, it will either
  return the answer
      \begin{verbatim}
        ERROR { explanation }
      \end{verbatim}
      if something has gone wrong, or the answer
      \begin{verbatim}
CATEGORY {category} DISTRIBUTION { category 1 } DISTANCE { 1.000000} NEIGHBORS
ENDNEIGHBORS
      \end{verbatim}
      where the presence of the {\tt DISTRIBUTION}, {\tt DISTANCE},
      {\tt MATCH\_DEPTH}, {\tt CONFIDENCE} and {\tt NEIGHBORS} parts
      depends upon the current verbosity settings ({\tt +vdb}, {\tt
        +vdi}, {\tt +vmd}, {\tt +vcf} and {\tt +vn}, respectively).
      Note that if the last string on the answer line is {\tt
        NEIGHBORS}, the server will proceed to return lines of nearest
      neighbor information until it prints the keyword {\tt
        ENDNEIGHBORS}.
\item {\tt set option}\\ where option is specified as a string of
      commandline options. Only the following commandline
      options are valid in this context: {\tt k, m, d, B, L, Q, w, v,
      x}. The setting of an option in this client does not affect the
      behavior of the server towards other clients. The server replies
      either with {\tt OK} or with {\tt ERROR \{explanation\}}.
    \item {\tt query}\\
      queries the server for a list of current settings. Returns a
      number of lines with status information, starting with a line
      that says {\tt STATUS}, and ending with a line that says {\tt
        ENDSTATUS}. For example:

\begin{footnotesize}
\begin{verbatim}
STATUS
FLENGTH              : 0
MAXBESTS             : 500
NULL_VALUE           :
TREE_ORDER           : G/V
DECAY                : Z
INPUTFORMAT          : Column
SEED                 : -1
DECAYPARAM           : 1.000000
SAMPLE_WEIGHTS       : -
IGNORE_SAMPLES       : +
PROBALISTIC          : -
VERBOSITY            : F
EXACT_MATCH          : -
USE_INVERTED         : -
GLOBAL_METRIC        : Overlap
METRICS              :
NEIGHBORS            : 1
PROGRESS             : 100000
TRIBL_OFFSET         : 0
IB2_OFFSET           : 0
WEIGHTING            : GRW
ENDSTATUS
\end{verbatim}
\end{footnotesize}

\item {\tt exit}\\
  closes the connection between this client and the server.
\end{description}

\section{The HTTP protocol}
\label{httpproto}

With {\tt protocol=http}, the server answers HTTP {\tt GET} requests,
so that a trained TiMBL base can be queried directly from a web
browser or with tools such as {\tt curl}. The result is an XML
document.

The path of the request selects the base, and the query string
contains one or more actions of the form {\tt attribute=value},
separated by \&{}-signs:

\begin{description}
\item {\tt classify="testcase"} : classify {\tt testcase}. The
  (URL-encoded) instance may be enclosed in single or double quotes.
  This action may be repeated to classify several instances in one
  request.
\item {\tt set=option} : set a TiMBL option for this request, as with
  the {\tt set} command of the TCP protocol; a leading {\tt -} may be
  omitted (e.g.\ {\tt set=k 3}).
\item {\tt show=settings} or {\tt show=weights} : include the current
  settings, or the current feature weights, in the XML answer.
\end{description}

For example, with the configfile of Section~\ref{configfile} running
on {\tt localhost:7000} with {\tt protocol=http}:

\begin{footnotesize}
\begin{verbatim}
> curl 'http://localhost:7000/dimin0?classify="=,=,=,=,=,=,=,=,+,p,e,=,?"'
\end{verbatim}
\end{footnotesize}

returns an XML document along the lines of

\begin{footnotesize}
\begin{verbatim}
<?xml version="1.0"?>
<TiMblResult algorithm="IB1">
  <classification>
    <input>=,=,=,=,=,=,=,=,+,p,e,=,?</input>
    <category>T</category>
    <distribution>{ T 1.00000 }</distribution>
    <distance>0</distance>
  </classification>
</TiMblResult>
\end{verbatim}
\end{footnotesize}

where, as in the TCP protocol, the presence of the {\tt distribution},
{\tt distance}, {\tt confidence} and {\tt match\_depth} elements
depends on the verbosity settings of the base.

\section{The JSON protocol}
\label{jsonproto}

Since version 1.13, TimblServer can communicate in JSON. With {\tt
  protocol=json}, client and server exchange newline-delimited JSON
objects over the socket: every request is a single JSON object on one
line, and every answer is a single JSON object (or array) on one
line. {\bf Note:} the JSON syntax might still change in future
versions, so handle with care.

On connecting, the server sends a greeting. When the server hosts a
single (default) base, the greeting is:

\begin{verbatim}
{"status":"ok"}
\end{verbatim}

When several bases are available, the greeting lists them, and a base
must be selected before classifying:

\begin{verbatim}
{"status":"ok","available_bases":["dimin0","dimin1","dimin2"]}
\end{verbatim}

Every request object has a {\tt "command"} key, and optionally a {\tt
  "param"} key (a single string) or a {\tt "params"} key (an array of
strings). The commands mirror those of the TCP protocol:

\begin{description}
\item {\tt base}\\
  select a base, e.g.\ \verb|{"command":"base","param":"dimin0"}|.\\
  The server answers \verb|{"base":"dimin0"}|, or an error object.
\item {\tt classify}\\
  classify one instance given as {\tt "param"}, or several instances
  given as {\tt "params"}:
\begin{footnotesize}
\begin{verbatim}
{"command":"classify","param":"=,=,=,=,=,=,=,=,+,p,e,=,?"}
{"command":"classify","params":["-,=,=,=,+,k,e,=,?","=,=,=,=,+,p,e,=,?"]}
\end{verbatim}
\end{footnotesize}
  The answer for a single instance is a JSON object; for multiple
  instances it is a JSON array of such objects. Each result object has
  a {\tt "category"} key, and, depending on the verbosity settings of
  the base ({\tt +vdi}, {\tt +vdb}, {\tt +vmd}, {\tt +vcf}, {\tt
    +vn}), the additional keys {\tt "distance"}, {\tt "distribution"},
  {\tt "match\_depth"}, {\tt "confidence"} and {\tt "neighbors"}:
\begin{footnotesize}
\begin{verbatim}
{"category":"T","distance":0.0,"distribution":"{ T 1.00000 }"}
\end{verbatim}
\end{footnotesize}
\item {\tt set}\\
  set TiMBL options for this session, e.g.\
  \verb|{"command":"set","param":"-k 3"}|.  The server answers
  \verb|{"status":"ok"}| or an error object.
\item {\tt query} (or {\tt show})\\
  with {\tt "param"} {\tt "settings"} or {\tt "weights"}, return the
  current settings or the current feature weights of the selected base
  as a JSON object, e.g.\
  \verb|{"command":"query","param":"weights"}| returns
\begin{footnotesize}
\begin{verbatim}
{"weighting":"gr","weights":[0.024891,0.021873,...,0.643681]}
\end{verbatim}
\end{footnotesize}
\item {\tt exit}\\
  close the session. The server answers \verb|{"status":"closed"}|.
\end{description}

Whenever a request cannot be handled, the server answers with an error
object:

\begin{verbatim}
{"status":"error","message":"Unknown command: 'clasify'"}
\end{verbatim}

A complete session with the configfile of Section~\ref{configfile}
(with {\tt protocol=json}) looks like this, with lines sent by the
client marked {\tt >}:

\begin{footnotesize}
\begin{verbatim}
  {"status":"ok","available_bases":["dimin0","dimin1","dimin2"]}
> {"command":"base","param":"dimin0"}
  {"base":"dimin0"}
> {"command":"classify","param":"=,=,=,=,=,=,=,=,+,p,e,=,?"}
  {"category":"T","distance":0.0,"distribution":"{ T 1.00000 }"}
> {"command":"exit"}
  {"status":"closed"}
\end{verbatim}
\end{footnotesize}

\chapter{The timblclient program}
\label{client}

The distribution includes {\tt timblclient}, a simple client program
for the TCP protocol, which can be used to test a running server or to
classify a test file against it:

{\tt > timblclient -h host -p port [-i inputfile] [-o outputfile]
  [--batch] [-b basename]}

\begin{description}
\item {\tt -h <host>} : connect to the server on 'host'.
\item {\tt -p <port>} : connect to 'port' on 'host'.
\item {\tt -i <inputfile>} : a series of timblserver commands, as
  described in Section~\ref{tcpproto}. Default is standard input.
\item {\tt -o <outputfile>} : where to put the results. Default is
  standard output.
\item {\tt --batch} : when {\tt --batch} is specified, the input from
  'inputfile' is not interpreted as timblserver commands, but as a
  'normal' Timbl test file: every line is sent to the server 'as is',
  prepended by a {\tt classify} instruction.
\item {\tt -b <basename>} : select 'basename' on the server before
  testing.
\end{description}

For example, classifying a complete test file against the server from
the Quick-start Tutorial (Chapter~\ref{tutorial}):

{\tt > timblclient -h localhost -p 7000 --batch -i dimin.test -o dimin.out}

\clearpage

\begin{thebibliography}{1}

\bibitem{Daelemans+26}
Walter Daelemans, Jakub Zavrel, Ko van der Sloot, and Antal van den
  Bosch (2026).
\newblock {\em {TiMBL}: Tilburg Memory-Based Learner, version 7.0,
  Reference Guide}.
\newblock Available from \timblmanualurl.

\bibitem{Daelemans+05}
Walter Daelemans and Antal van den Bosch (2005).
\newblock {\em Memory-Based Language Processing}.
\newblock Cambridge University Press, Cambridge, UK.

\end{thebibliography}

\end{document}
